\section{Experiments and Results}
\label{sec:experiments}

\subsection{Experimental Setup}

\textbf{Dataset.} The VQA-RAD dataset \cite{lau2018dataset} contains 315 radiology images across multiple modalities (X-ray, CT, MRI) with 3,515 clinician-generated QA pairs. The standard train/test split provided by the dataset authors is used.

\textbf{Hardware.} All experiments are conducted on NVIDIA Tesla T4 $\times$ 2 GPUs with 14.6 GB VRAM each. Training is performed in mixed precision (FP16).

\textbf{Evaluation Metrics.} Closed-ended questions are evaluated using Exact Match (EM) accuracy. Open-ended questions are evaluated using BLEU-1 through BLEU-4 and ROUGE-L F1 scores.

\subsection{Results}

% ===== PLACEHOLDER: fill in after running experiments =====
\begin{table}[htbp]
\caption{Evaluation results on the VQA-RAD test set.}
\label{tab:results}
\centering
\begin{tabular}{lcc}
\toprule
\textbf{Metric} & \textbf{QLoRA (ours)} & \textbf{Zero-shot} \\
\midrule
Closed EM (\%)       & 68.13 & - \\
Open BLEU-1          & 0.3388 & - \\
Open BLEU-4          & 0.0809 & - \\
Open ROUGE-L         & 0.3774 & - \\
\midrule
Trainable params     & 19.1M & N/A \\
Training time        & $\sim$2.5 hours & N/A \\
Peak VRAM (GB)       & 14.6 & - \\
\bottomrule
\end{tabular}
\end{table}
% ===== END PLACEHOLDER =====

Table~\ref{tab:results} presents the comparison between the QLoRA fine-tuned model and the zero-shot baseline (LLaVA-1.5-7B without fine-tuning). The results indicate that QLoRA allows the LLaVA-1.5-7B model to adapt effectively to the medical domain. With less than 1\% of parameters trained, the model achieved a 68.13\% Exact Match on closed-ended questions and a 0.3388 BLEU-1 score on open-ended questions. Furthermore, the peak VRAM footprint was successfully contained within the 14.6 GB limit of a single Tesla T4 GPU (using model pipeline parallelization across two GPUs), proving the efficiency of 4-bit quantization combined with LoRA.

% ===== PLACEHOLDER: add loss curve figure after training =====
% \begin{figure}[htbp]
% \centering
% \includegraphics[width=\columnwidth]{figures/loss_curve.png}
% \caption{Training and validation loss curves over 5 epochs.}
% \label{fig:loss}
% \end{figure}
% ===== END PLACEHOLDER =====

\subsection{Analysis}

Qualitative analysis suggests the model is highly capable of answering binary "yes/no" questions (closed-ended) regarding the presence of abnormalities. However, for open-ended questions, generating precise medical terminology remains challenging, as reflected by the lower BLEU-4 score (0.0809). The Exact Match score highlights strong structural understanding, while the generative aspects indicate room for improvement, possibly by increasing the LoRA rank or expanding the fine-tuning dataset.
